\documentclass[10pt,conference]{IEEEtran}
\usepackage{cite}
\usepackage{amsmath,amssymb}
\usepackage{booktabs}
\usepackage{array}
\usepackage{url}
\usepackage[T1]{fontenc}
\renewcommand{\arraystretch}{1.08}
\begin{document}

\title{Can We Stop Ourselves from Aging?\\
A Design-Only Framework for Testing Bounded Changes in Human Aging Trajectories}

\author{Survey--Idea--ExperimentDesign--Author Research Workflow}

\maketitle

\begin{abstract}
Whether people can ``stop aging'' is not a single biomedical endpoint. Chronological time, biological change, functional decline, disease incidence, and mortality are related but nonidentical quantities. A late-life plateau in a conditional mortality hazard, where supported by a specified demographic data set, neither proves that an individual has stopped biological aging nor establishes an intervention. Likewise, an intervention-associated shift in one epigenetic, metabolic, inflammatory, or physiologic measure is not a validated proof of health-span extension. This design-only paper develops the Systematic Trial Architecture for Biological Longevity and Aging Evaluation (STABLE-AGE), an auditable protocol for assessing a candidate geroscience intervention in a declared population and time horizon. STABLE-AGE requires a frozen endpoint hierarchy spanning safety, calibrated multidomain biomarkers, function, and clinical-event follow-up; a randomized comparative core where ethical and feasible; blinded assay and outcome procedures; explicit missing-data and survival-selection analyses; held-out quality tests; and external replication. It treats molecular hallmarks as mechanism hypotheses, not a clinical diagnosis, and requires concordance across domains without allowing average biomarker improvement to compensate for a safety breach. The proposal reports no participant enrollment, treatment exposure, laboratory result, functional outcome, adverse event, mortality effect, or efficacy estimate. Its conclusion is conditional: no current evidence establishes that human aging can be stopped, but carefully bounded causal claims about specified aging-related trajectories can be made progressively more informative through reproducible, safety-centered studies.
\end{abstract}

\begin{IEEEkeywords}
aging biology, geroscience, biological age, biomarkers, health span, randomized trials, senescence, clinical translation, causal inference.
\end{IEEEkeywords}

\section{Introduction}
\label{sec:introduction}

The wish to stop aging compresses several scientific questions into one phrase. It may refer to avoiding a particular age-related disease, maintaining mobility and cognition, extending life, lowering a mortality hazard, repairing a cellular defect, changing a biomarker, or preventing every future biological change. These are not interchangeable endpoints. A clock can advance while a person remains functional; a clock can change after an exposure without demonstrating a durable health benefit; a clinical event can be delayed by disease-specific care without modifying the broad biology of aging. The first requirement for research in this area is therefore a precise claim boundary.

This paper uses \emph{aging-related trajectory} to denote a preregistered, multidomain pattern of change measured over a declared horizon. It does not claim that a single latent variable completely describes human aging. Rather, it recognizes that molecular damage, nutrient sensing, immune function, tissue repair, metabolic state, cognition, mobility, disease susceptibility, and mortality can diverge. A valid study must make such divergence visible. It must not choose a favorable readout after treatment, call a demographic pattern a biological mechanism, or call a mechanism result a clinical cure.

The biological literature provides a productive but incomplete map. The original hallmarks of aging organized processes including genomic instability, telomere attrition, epigenetic alteration, loss of proteostasis, deregulated nutrient sensing, mitochondrial dysfunction, cellular senescence, stem-cell exhaustion, and altered intercellular communication~\cite{lopez2013}. The expanded framework adds or elevates disabled macroautophagy, chronic inflammation, and dysbiosis~\cite{lopez2023}. These categories support hypotheses about how an intervention might act and where harm might emerge. They do not define a universal laboratory test for whether a person is ``aging,'' and they do not establish that manipulating one category will safely improve all others.

The question has also been linked to exceptional-age demography. Barbi \emph{et al.}\ reported a late-life mortality plateau in a carefully constructed population of longevity pioneers~\cite{barbi2018}. This observation deserves careful interpretation. A mortality hazard at age $a$, conditional on having survived to age $a$, is a population quantity; it is not a measurement of damage accumulation inside every survivor. Errors in age validation, population selection, frailty heterogeneity, migration, model choices, sparse event counts, and confidence intervals can alter an apparent late-age pattern. Even if a plateau were accepted for a specified population, it would not mean that all molecular or functional change has ended, nor that a tested intervention created the plateau.

The contribution of this report is STABLE-AGE: the \emph{Systematic Trial Architecture for Biological Longevity and Aging Evaluation}. It creates a design-only route from mechanism hypothesis to bounded inference. The framework requires researchers to declare the population, intervention class, comparator, follow-up window, safety boundary, biological assays, functional outcomes, clinical-event policy, missing-data model, and replication threshold before examining outcomes. It does not propose a treatment regimen or give medical advice. It reports no new experiment, human result, clinical benefit, or proof that aging can be stopped.

\section{Survey: What Existing Evidence Can and Cannot Establish}
\label{sec:survey}

\subsection{Chronological age, biological change, and conditional mortality}

Chronological age is measured without a biological assay; it is time since birth. Biological aging is an explanatory concept for age-associated change, and its proposed measures must be validated for a stated use. Mortality is an outcome caused by many proximal pathways, including disease, injury, infection, environmental exposure, care access, and biological vulnerability. Conflating these quantities invites false claims. A lower mortality rate in one cohort does not reveal a single cause. A favorable cellular readout does not guarantee a lower hazard. A person who reaches an exceptional age is not a control experiment showing that aging has vanished.

For a population with covariates $X$ and survival indicator $S$, the relevant late-age quantity is a conditional hazard,
\begin{align}
h(a\mid X,S(a)=1)&=\lim_{\Delta a\rightarrow 0} \nonumber\\
&\frac{\Pr[a\leq A<a+\Delta a\mid A\geq a,X]}{\Delta a},
\label{eq:conditional-hazard}
\end{align}
where $A$ is age at death and $a$ is attained age. Equation~\eqref{eq:conditional-hazard} makes the survivor conditioning explicit. A comparison of $h(105\mid X,S(105)=1)$ and $h(106\mid X,S(106)=1)$ concerns successive selected survivor populations. It cannot by itself identify the within-person rate of change of molecular damage, functional capacity, or a latent biological-age construct.

There are several reasons to maintain this distinction. First, exceptional ages are rare, so a result can be sensitive to record verification and model specification. Second, unobserved heterogeneity means that individuals with different vulnerabilities are not equally represented among later survivors. Third, an age-specific death rate is not an intervention endpoint unless the intervention, population, exposure duration, confounding structure, and competing risks are identified. STABLE-AGE separates an S0 demographic reanalysis from an S2 intervention experiment for this reason.

\subsection{Mechanism frameworks are maps, not efficacy certificates}

The hallmarks literature is valuable because it resists a single-cause account. An intervention affecting nutrient sensing, for example, may also affect immune function, metabolism, wound healing, muscle, cancer biology, or medication interactions. A senescence-targeting strategy may have distinct effects across cell types and disease contexts. The systems nature of the framework implies that a favorable response in one pathway should be accompanied by measurement of plausible countervailing pathways and function.

This is especially important for translation. Molecular processes have different time scales and tissue distributions. Blood can be practical to sample but may not represent brain, muscle, liver, bone marrow, or a disease-relevant tissue. A transient inflammatory response may be harmful, beneficial, or simply a context-dependent response. An assay may be analytically repeatable but biologically nonspecific. These properties do not make the assay useless; they define the interpretation it can support.

Table~\ref{tab:claim-boundary} distinguishes evidence classes. It deliberately refuses a shortcut from mechanism to clinical claim. A cell or animal result can motivate target selection. A human biomarker result can support target engagement. A randomized human functional outcome can support a bounded comparative claim. None, individually, establishes that aging as a whole has stopped.

\begin{table*}[!t]
\caption{Evidence boundary for claims about stopping aging. Each class may be scientifically useful, but it supports only the stated conclusion.}
\label{tab:claim-boundary}
\centering
\footnotesize
\begin{tabular}{@{}p{0.18\textwidth}p{0.25\textwidth}p{0.25\textwidth}p{0.24\textwidth}@{}}
\toprule
Evidence class & What it can establish & Critical limitations & Claim that remains unsupported \\
\midrule
Verified late-age demographic data & A conditional hazard pattern for the specified population, ages, model, and data-quality rules. & Survival selection, sparse events, measurement error, heterogeneity, and no intervention contrast. & Individuals have stopped biological aging, or an intervention has been validated. \\
Mechanistic cell or animal study & A hypothesis about a pathway, exposure, or tissue-specific mechanism under stated conditions. & Species, dose, tissue, environment, and outcome transfer may fail in humans. & Human health-span, safety, or lifespan benefit. \\
Human biomarker change & Analytical response or target engagement for the measured specimen and assay. & Surrogacy, assay drift, tissue specificity, regression to mean, and multiplicity. & Slowed aging or reduced disease risk without functional and clinical concordance. \\
Randomized function/clinical comparison & A bounded causal contrast for enrolled participants, comparator, endpoints, and follow-up. & Adherence, missingness, external validity, competing risks, and limited horizon. & Permanent, universal, or risk-free cessation of aging. \\
Replication and long follow-up & Reproducibility and more credible durability for the replicated boundary. & Still finite, population-specific, and unable to prove eternal survival. & Immortality or unrestricted longevity. \\
\bottomrule
\end{tabular}
\end{table*}

\subsection{Biomarkers and the surrogate problem}

Epigenetic clocks, proteomic panels, metabolomic signatures, inflammatory indices, imaging measures, physiologic tests, and composite scores each answer different measurement questions. Horvath's DNA-methylation age work illustrates how chronological-age-associated information can be extracted across tissues~\cite{horvath2013}. Levine \emph{et al.}\ developed a biomarker intended to link methylation patterns to lifespan and health-span-related phenotypes~\cite{levine2018}. Such work is important for measurement science. It does not make every clock interchangeable or automatically validate its change as a surrogate for an intervention's clinical benefit.

Surrogacy has a causal requirement. A biomarker must do more than correlate with age or future risk; its change under intervention must reliably convey the intervention's effect on outcomes of interest within the intended population and setting. A treatment can improve a biomarker while producing no functional benefit, or even net harm, through an off-target pathway. Conversely, a functional benefit can occur without a large change in the chosen biomarker. STABLE-AGE consequently makes all component domains visible and treats discordance as a result requiring explanation rather than as a reason to select the most favorable assay.

Let $Y_{ikt}$ be an observed value for participant $i$, domain $k$, and visit $t$. A useful measurement decomposition is
\begin{equation}
Y_{ikt}=\theta_{ik}(t)+b_{k,\mathrm{site}}+b_{k,\mathrm{batch}}+\varepsilon_{ikt},
\label{eq:measurement-decomposition}
\end{equation}
where $\theta_{ik}(t)$ is the underlying domain state, $b_{k,\mathrm{site}}$ and $b_{k,\mathrm{batch}}$ are premeasured site and batch effects, and $\varepsilon_{ikt}$ is residual measurement variation. Equation~\eqref{eq:measurement-decomposition} is not a proof that a latent state is identified; it states the sources of variation that must be calibrated before a biological interpretation is made.

\subsection{Intervention evidence and the health-span boundary}

Human intervention studies supply candidate evidence, not a present solution to aging. Calorie restriction has been examined in CALERIE for cardiometabolic risk and other intermediate outcomes~\cite{kraus2019}; a later analysis reported a slowed pace-of-aging measure in healthy adults under the trial's conditions~\cite{belsky2023}. These studies make an important methodological contribution because they expose intervention assignment, adherence, measured outcomes, and time horizon. They do not show that all forms of aging have stopped, and their eligibility and burden conditions matter for transportability.

Nutrient-sensing modulation is another active area. Mannick \emph{et al.}\ reported immune-function-related findings following mTOR inhibition in older adults~\cite{mannick2014}. This supports continued hypothesis-driven clinical investigation, not an inference that mTOR modulation is a broadly safe anti-aging treatment. Senescence-targeting work provides a complementary example: Justice \emph{et al.}\ reported an early, open-label pilot in idiopathic pulmonary fibrosis~\cite{justice2019}. A pilot can establish feasibility, tolerability signals, and endpoint learning, but it cannot settle comparative clinical efficacy or general biological aging.

Public-health evidence and clinical geroscience must also be separated. The National Institute on Aging summarizes established actions associated with healthy aging and the maintenance of function~\cite{nia2026}. These are valuable health practices, but they are not a license to label a behavior or product as having stopped aging. The research question in this report is whether a future, specified candidate intervention satisfies a demanding causal and safety standard. It should not be read as a recommendation to start, stop, substitute, or combine any medication, diet, supplement, or procedure.

\section{Idea: STABLE-AGE as a Bounded Causal Claim}
\label{sec:idea}

STABLE-AGE is a study architecture rather than a therapy. Its core question is: \emph{for a declared eligible population, does intervention class $a$ change a preregistered multidomain aging-related trajectory compared with a defined control during a stated horizon, with compatible functional or clinical evidence and without an unacceptable safety signal?} The question is deliberately narrower than ``Can we stop aging?'' because it can be falsified by an endpoint, a safety event, failed assay calibration, absent concordance, or nonreplication.

The architecture has five principles. First, define the claimed object before measuring it. Second, distinguish technical assay validity from biological interpretation and from clinical surrogacy. Third, make health and safety hard constraints rather than weights that can be offset by a favorable average biomarker. Fourth, protect causal inference with comparison, blinding where feasible, prespecified analysis, and transparent missingness handling. Fifth, label every conclusion by its evidence layer: observation, measurement inference, causal estimate, or long-horizon extrapolation.

For calibrated domain summaries $Z_k(t)$ with locked directions and scaling, STABLE-AGE may define a transparent composite,
\begin{equation}
B(t)=\sum_{k=1}^{K}w_k Z_k(t),\qquad \sum_{k=1}^{K}w_k=1,
\label{eq:multidomain-composite}
\end{equation}
where $w_k$ is prespecified from external validation or a protocol rule and $K$ is the number of retained domains. Equation~\eqref{eq:multidomain-composite} is a reporting device, not a claim that all domains share a single biological cause. The full vector $(Z_1,\ldots,Z_K)$ must accompany $B(t)$ so that a benefit in one domain cannot conceal deterioration in another.

The corresponding within-person change rate over a declared horizon $T$ is
\begin{equation}
r_i(T)=\frac{B_i(T)-B_i(0)}{T},
\label{eq:trajectory-rate}
\end{equation}
where $B_i(0)$ and $B_i(T)$ use the same locked assay and quality-control rules. Equation~\eqref{eq:trajectory-rate} permits a bounded comparison of rates; it does not establish that $r_i(T)=0$ in an unmeasured future or that all relevant biological changes share the same rate.

STABLE-AGE therefore defines four hypotheses. Under $H_0$, an apparent result is attributable to chance, imbalance, batch drift, selective retention, behavior change, regression to the mean, or survival selection. Under $H_1$, treatment assignment changes the preregistered trajectory and is concordant with independently measured function or clinical evidence in the enrolled population. Under $H_2$, a biomarker moves without such concordance, which supports target engagement or measurement research only. Under $H_3$, effects are heterogeneous, context-dependent, or offset by harm; broad anti-aging language is then prohibited.

\begin{table*}[!t]
\caption{STABLE-AGE endpoint hierarchy. The safety and quality gates are noncompensatory.}
\label{tab:endpoint-hierarchy}
\centering
\footnotesize
\begin{tabular}{@{}p{0.17\textwidth}p{0.28\textwidth}p{0.25\textwidth}p{0.22\textwidth}@{}}
\toprule
Tier & Examples of preregistered content & Main inferential role & Prohibited shortcut \\
\midrule
Safety hard gates & Serious adverse events, clinically meaningful laboratory changes, deterioration in function, intervention-specific risks, and stop rules. & Determines whether benefit claims are eligible for interpretation. & Offsetting a safety failure with a favorable composite score. \\
Assay and quality gates & Sample identity, technical replicates, blinded duplicates, batch/site controls, calibration acceptance, and drift monitoring. & Determines whether observed change is measurement-credible. & Calling a batch effect biological target engagement. \\
Primary multidomain trajectory & Locked molecular, physiologic, immune, and functional-domain summaries with complete component reporting. & Tests the bounded trajectory hypothesis. & Replacing the component vector with a single favorable clock. \\
Concordance outcomes & Independent functional measures and, when feasible, blinded clinical adjudication. & Tests whether the trajectory has compatible consequences. & Treating correlation with age as clinical validation. \\
Long-horizon event extension & Incident multimorbidity, hospitalization, disability, and mortality with follow-up completeness. & Describes durability and event patterns. & Inferring mortality benefit from a short biomarker study. \\
\bottomrule
\end{tabular}
\end{table*}

\section{ExperimentDesign: Integrated Human-Research Protocol}
\label{sec:design}

\subsection{Scenario ladder and study boundary}

The protocol progresses through scenario classes rather than claiming that one study can settle a universal question. S0 is a demographic reanalysis used to keep exceptional-age hazard models separate from intervention efficacy. S1 is an assay-calibration program. S2 is a randomized comparative core where ethical and feasible. S3 is a longer observational event extension, which retains association language unless its causal assumptions are met. S4 is an external replication in an independent site or cohort. The ladder prevents an investigator from promoting an early mechanistic or measurement result directly to a public claim of stopped aging.

Table~\ref{tab:scenario-ladder} specifies the evidence each scenario may provide. The boundary includes population eligibility, baseline phenotype, comorbidities, medication stability, relevant reproductive considerations, frailty limits, access to usual care, specimen schedule, intervention exposure, comparator, follow-up policy, and endpoints. These choices must be documented before recruitment. They are scientific assumptions, not administrative details, because they determine the population to which an inference applies.

\begin{table*}[!t]
\caption{STABLE-AGE scenario ladder. A later scenario adds evidence but does not erase limits of a finite study.}
\label{tab:scenario-ladder}
\centering
\footnotesize
\begin{tabular}{@{}p{0.14\textwidth}p{0.30\textwidth}p{0.28\textwidth}p{0.20\textwidth}@{}}
\toprule
Scenario & Required evidence and controls & Failure disqualifier & Permissible conclusion \\
\midrule
S0: demographic reanalysis & Verified age records, inclusion logic, censoring policy, competing-risk model, and sensitivity analyses. & Apparent plateau depends on unexamined age errors, population definition, or model choice. & A conditional demographic pattern is described for the stated data boundary. \\
S1: assay calibration & Technical replicates, blinded duplicates, sample swaps, batch/site controls, and predeclared quality thresholds. & Identity, repeatability, drift, or cross-platform agreement is inadequate. & The assay is qualified for a stated measurement purpose. \\
S2: randomized comparative core & Concealed allocation, comparator, outcome schedule, adherence records, blinded assays and adjudication where feasible. & Safety breach, protocol instability, unblinding bias, or endpoint switching. & A bounded causal contrast is estimated for enrolled participants. \\
S3: event extension & Complete event definitions, follow-up provenance, loss-to-follow-up report, and competing-risk analysis. & Informative attrition or care changes are ignored. & Longer-horizon event association or protocol-specified extension evidence is reported. \\
S4: external replication & Frozen analysis, independent staff/site, prespecified population differences, and full null-result reporting. & Interpretation depends on retuning analysis after seeing replication data. & Reproducibility support for the replicated boundary is obtained. \\
\bottomrule
\end{tabular}
\end{table*}

The report does not select a dose, product, supplement, off-label medicine, or invasive procedure. Such selection would require candidate-specific pharmacology, toxicology, interaction, manufacturing, clinical, ethics, and regulatory review. STABLE-AGE instead accepts a versioned intervention description $a$, a matched control $c$, and a protocol stating what concomitant care, rescue treatment, adherence support, and discontinuation conditions are permitted. This preserves clinical equipoise and prevents the general framework from becoming disguised treatment advice.

\subsection{Sampling, randomization, and endpoint freezing}

At baseline, each participant receives a versioned phenotype record that includes the protocol-defined eligibility measures, clinical history, medication record, exposure context, specimen timing, and functional assessment. Randomization, when used, is stratified on only the predeclared factors that affect intervention safety, assay interpretation, or expected clinical risk. The allocation procedure, concealment method, unblinding conditions, visit window, specimen logistics, laboratory assignment, and data lock date are maintained in a protocol register.

Endpoint freezing protects against flexible interpretation. Before unblinding, the protocol fixes domain direction, transformations, weights, quality thresholds, visit windows, primary contrast, covariate set, missing-data primary model, multiplicity family, and sensitivity analyses. It also states which measures are confirmatory and which are exploratory. No biomarker, subgroup, time window, or composite weight may be substituted merely because it provides a more attractive result. The purpose is not to suppress discovery; exploratory findings may be reported as such, with their full search context.

The primary mixed model can be expressed as
\begin{align}
Y_{ikt}={}&\beta_{0k}+\beta_{1k}t+\beta_{2k}A_i+\beta_{3k}(A_i t) \nonumber\\
&+\gamma_k^{\mathsf{T}}X_i+u_{ik}+v_{\mathrm{site}(i),k}+\varepsilon_{ikt},
\label{eq:mixed-model}
\end{align}
where $A_i$ denotes assignment, $X_i$ is the locked baseline covariate vector, $u_{ik}$ is a participant-level random effect, $v_{\mathrm{site}(i),k}$ is a site-level effect, and $\beta_{3k}$ is the prespecified differential change parameter for domain $k$. Equation~\eqref{eq:mixed-model} is illustrative; the final model family must match the scale and distribution of each endpoint. Its key feature is that treatment-by-time is prespecified rather than inferred from a selected before-after comparison.

For the locked composite, the principal contrast is
\begin{equation}
\Delta_B(T)=\mathbb{E}[r_i(T)\mid A_i=1]-\mathbb{E}[r_i(T)\mid A_i=0],
\label{eq:composite-contrast}
\end{equation}
with $r_i(T)$ defined in Equation~\eqref{eq:trajectory-rate}. Equation~\eqref{eq:composite-contrast} is reported with its uncertainty, the component-domain estimates, and all safety outcomes. It does not make the composite a replacement for function or clinical events.

\subsection{Safety, function, and clinical-event follow-up}

Safety is a gate, not a utility weight. The study safety plan must include adverse-event capture, clinically meaningful laboratory monitoring, intervention-specific surveillance, functional deterioration checks, independent oversight, protocol-defined pause or stopping conditions, and care-referral procedures. The choice of monitoring depends on the candidate intervention and must be evaluated by qualified clinicians, ethics boards, and regulators. Participants remain entitled to usual clinical care; research measurements do not substitute for diagnosis or management.

Function is measured through a preregistered set of outcomes suited to the enrolled population and capable of blinded or standardized collection when feasible. A function signal that is compatible with a biomarker direction is more informative than a biomarker alone, but it remains finite and population-specific. Clinical events are retained in the protocol even when the primary horizon cannot credibly estimate event differences. This prevents a short study from implying disease prevention simply because the intended clinical endpoint was omitted.

The safety-and-concordance gate can be stated as a logical rule,
\begin{equation}
G=\mathbb{I}_{\mathrm{safety}}\mathbb{I}_{\mathrm{quality}}\mathbb{I}_{\mathrm{trajectory}}\mathbb{I}_{\mathrm{concordance}}\mathbb{I}_{\mathrm{replication}},
\label{eq:noncompensatory-gate}
\end{equation}
where each indicator equals one only if its preregistered criterion is met. Equation~\eqref{eq:noncompensatory-gate} means that $G=0$ whenever a hard safety gate, assay-quality criterion, primary trajectory rule, concordance rule, or replication requirement fails. It intentionally prevents a high average biological score from compensating for a serious adverse outcome or a failed replication.

\subsection{Missing data, adherence, and survivor bias}

Older populations and long studies can have hospitalization, changing care, frailty, inability to travel, withdrawal, death, and assay failure. These processes may be affected by the intervention or by the outcome, so deleting incomplete records can create a false benefit or obscure harm. STABLE-AGE requires visit provenance, reason-for-missingness coding, discontinuation dates, intervention adherence, rescue care, assay failure logs, and event follow-up for all possible participants.

One sensitivity approach uses stabilized inverse-probability weights. Let $L_i(s)$ denote the locked baseline, outcome-history, health, and adherence covariates that predict retention at visit $s$. Then
\begin{equation}
W_i(t)=\prod_{s\leq t}\frac{q_s(1\mid A_i)}{q_s(1\mid A_i,L_i(s))},
\label{eq:retention-weight}
\end{equation}
where $q_s$ is the specified conditional probability of observed retention. Equation~\eqref{eq:retention-weight} is a model-based sensitivity tool, not a guarantee that missingness is solved. Its assumptions, weight diagnostics, alternative models, and departures from positivity must be reported.

Survival selection receives a separate analysis. A participant's absence after death cannot be imputed as an ordinary biomarker value. Event analyses must state the estimand, competing-risk handling, censoring rule, and relation between survival and repeated measurement. The demographic S0 analysis also cannot be used to infer the causal effect of S2 treatment assignment. These separations protect against the attractive but incorrect inference that observed later survivors demonstrate treatment-induced arrest of aging.

\section{Systematics, Calibration, and Decision Rules}
\label{sec:systematics}

Before any participant-level result is interpreted, STABLE-AGE executes a calibration and fault-injection campaign. Technical replicates establish repeatability. Blinded duplicate specimens test operator and laboratory agreement. Sample swaps test identity controls. Batch and site controls test whether a laboratory transition can imitate a biological change. Null analyses test whether a pipeline creates a treatment effect when assignment is absent. Synthetic injections add a known batch shift, temporary acute illness, nonadherence, selective dropout, or delayed follow-up so the analysis can be tested for false reassurance.

The point of these tests is not to create synthetic evidence of efficacy. They test the honesty of the measurement and analysis pipeline. If a known batch shift is reported as target engagement, if a null allocation produces a treatment benefit, or if selective dropout is hidden by complete-case analysis, the system is not ready for a biological conclusion. Failed calibration is an informative result that narrows the claim boundary.

Multiplicity must also be explicit. For $m$ confirmatory hypotheses with ordered $p$-values $p_{(1)}\leq\cdots\leq p_{(m)}$, a simple family-wise step-down criterion compares
\begin{equation}
p_{(j)}\leq \frac{\alpha}{m-j+1}
\label{eq:holm-rule}
\end{equation}
at step $j$ under the prespecified procedure. Equation~\eqref{eq:holm-rule} is included to make the confirmatory family visible; an implementation must specify the hypotheses, alpha allocation, multiplicity method, and exploratory analyses before data lock. It is not a substitute for well-chosen endpoints or external replication.

Table~\ref{tab:decision-gates} gives the progression rule. Each gate closes a specific inferential dependency. This is more useful than a generic maturity score because it identifies whether the blocking uncertainty is safety, measurement, causal comparison, functional relevance, long-horizon events, or reproducibility. A failed gate does not invalidate all upstream evidence; it restricts the language that evidence may support.

\begin{table*}[!t]
\caption{Decision gates for a bounded claim about an aging-related trajectory. Passing a gate does not establish that aging has stopped.}
\label{tab:decision-gates}
\centering
\footnotesize
\begin{tabular}{@{}p{0.15\textwidth}p{0.27\textwidth}p{0.28\textwidth}p{0.22\textwidth}@{}}
\toprule
Gate & Required evidence & Failure disqualifier & Permissible conclusion \\
\midrule
D0: construct and boundary & Population, intervention, comparator, horizon, endpoint dictionary, safety plan, and intended claim are locked. & The phrase ``aging'' is undefined or the claim changes after data review. & A testable protocol exists. \\
D1: measurement credibility & Identity checks, technical replicates, duplicate specimens, calibration, drift assessment, and source logs pass. & Assay, site, batch, or sample effects are unresolved. & A stated assay is credible for its declared measurement purpose. \\
D2: comparative causal core & Allocation/comparator, adherence, missingness model, blinded procedures where feasible, and complete outcome provenance are preserved. & Endpoint switching, unaddressed informative attrition, or loss of protocol integrity. & A bounded assignment contrast may be estimated. \\
D3: health relevance & Functional or clinical measures are compatible with primary trajectory direction and every safety gate remains satisfied. & Biomarker-only result, functional deterioration, or clinically important safety concern. & Target engagement plus limited health relevance is supported for the stated horizon. \\
D4: durability and replication & Event extension, frozen analysis, independent replication, and full null/harm reporting are available. & Retuned replication analysis, failed replication, or concealed adverse outcomes. & A reproducible, condition-specific trajectory claim is supported. \\
\bottomrule
\end{tabular}
\end{table*}

\section{Reproducibility, Ethics, and Interpretation Safety}
\label{sec:reproducibility}

The reproducibility package includes the registered protocol, analysis plan, endpoint dictionary, intervention version, randomization method as permitted, consent and privacy constraints, assay specifications, sample-processing metadata, calibration logs, laboratory and batch identifiers, code, container or software versions, all exclusions, missingness diagnostics, adverse-event summaries, and the distinction between confirmatory and exploratory analyses. De-identified analysis-ready data should be released where participant consent, privacy, and governance permit. Where release is not possible, an access process and an independent reproducibility audit plan should be stated.

Every conclusion must also identify its evidence layer. An \emph{observation} is a measured value or event. A \emph{measurement inference} states what a calibrated assay represents within its scope. A \emph{causal estimate} relies on its comparison design and assumptions. A \emph{long-horizon extrapolation} combines the preceding layers with unobserved future conditions. The report cannot turn a short-duration biomarker observation into a lifetime claim without clearly crossing these layers.

Ethical interpretation is inseparable from design. Aging-related claims can be commercially and emotionally salient, especially for people with illness or fear of disability. A study should not imply that enrollment confers treatment, should disclose uncertainty and alternatives to potential participants, and should maintain independent safety oversight. Research communication must distinguish an intervention hypothesis from clinical availability. The NIA's public discussion of healthy aging emphasizes that many factors affect health and function with age~\cite{nia2026}; that public-health context should not be conflated with experimental proof for a candidate geroscience intervention.

Equity and external validity also require attention. A biomarker model calibrated in one ancestry, sex, age range, disease profile, socioeconomic context, or care setting may not transport to another. Recruitment and retention procedures may preferentially include healthier, more resourced people. Medication access and monitoring capacity can affect apparent safety. STABLE-AGE therefore reports population coverage, representation gaps, site context, data provenance, and the limits on generalization rather than treating a statistically adjusted estimate as universal.

\section{Threats to Validity and Research Priorities}
\label{sec:validity}

The most immediate validity threat is outcome substitution. A study may advertise a general anti-aging result after observing one selected clock or laboratory value. The countermeasure is the locked endpoint hierarchy and component reporting in Table~\ref{tab:endpoint-hierarchy}. A second threat is survivor bias, particularly when extreme-age demography is cited as biological evidence. The countermeasure is to preserve the conditioning in Equation~\eqref{eq:conditional-hazard}, audit records, and keep S0 separate from the treatment comparison.

A third threat is intervention complexity. A combination of diet, exercise, medication changes, monitoring, social support, and study attention can affect measured outcomes through several pathways. A global result may be valuable, but it cannot automatically be attributed to one biological target. The protocol must state the estimand: the effect of assignment to the complete strategy, the effect while adhering to it, or a specified component effect. Claims stronger than the design can identify are prohibited.

A fourth threat is time-scale mismatch. Mechanistic signals may change in days or weeks, whereas disability, multimorbidity, or mortality can require much longer observation and different sample structures. This does not make short studies futile. It makes their conclusion narrower: analytical calibration, target engagement, or a short-horizon trajectory comparison. Long-horizon outcomes require their own follow-up design, event definitions, retention plan, and uncertainty reporting.

Research priority should follow value of information. The most valuable questions are those whose answers change a decision gate: whether a biomarker is sufficiently repeatable and responsive; whether an intervention's plausible harms exceed its anticipated benefit; whether an observed change maps to function; whether retention depends on emerging health; whether the effect persists after intervention discontinuation; and whether an independent cohort reproduces the result. An incremental gain in a selected clock has limited research value if it cannot change safety, function, causal interpretation, or replication readiness.

\section{Discussion}
\label{sec:discussion}

Can people stop aging? Current science does not support an affirmative answer. The question becomes more useful when transformed from a promise into a set of falsifiable, ethically supervised measurements. An intervention can be investigated for its effect on a specified biological process, a multidomain trajectory, a functional endpoint, or a clinical event pattern. Each claim needs the population, exposure, comparator, outcome, duration, and uncertainty to be visible. None should be expanded into ``aging has stopped'' because a finite study has not observed all future physiological, environmental, and disease pathways.

This conclusion is not pessimistic. It identifies a productive scientific path. Hallmarks research can guide mechanism hypotheses. Biomarker science can improve measurement and trial efficiency. Randomized comparative studies can estimate bounded causal effects. Functional and clinical outcomes can protect against surrogate optimism. Event extensions can examine durability. Replication can test whether an apparent effect is robust. The ladder is demanding because public claims in this area are consequential, but it turns unknowns into experiments instead of leaving them as rhetorical possibilities.

The framework also clarifies how to describe exceptional-age demography. A carefully verified hazard plateau, if reproduced under stated assumptions, is an interesting observation about mortality among survivors. It should stimulate work on heterogeneity, resilience, record validation, and the relation between demographic risk and biological state. It should not be used to infer that a person's aging rate is zero, that health deterioration is absent, or that a therapy has been discovered. STABLE-AGE maintains this distinction by design rather than by a disclaimer after results are known.

For translation, the next credible advance is not a universal biological-age number or a dramatic longevity promise. It is a transparent study that can fail cleanly: a candidate intervention with a measured exposure, a comparator, calibrated assays, a frozen endpoint hierarchy, safety oversight, a plan for missingness and survivor selection, and a replication boundary. If its result is null, harmful, discordant, or nonreplicable, the research program learns which dependency remains open. If it passes the gates, its resulting claim remains finite but becomes more useful and more trustworthy.

\section{Conclusion}
\label{sec:conclusion}

This four-stage research workflow concludes that human aging cannot currently be claimed to be stopped. A late-life mortality plateau is not equivalent to cessation of biological change, and mechanism or biomarker results are not equivalent to health-span or lifespan proof. STABLE-AGE provides a design-only framework for a more scientific question: whether a declared intervention changes a preregistered, multidomain aging-related trajectory for a specified population and horizon while meeting noncompensatory safety, measurement, functional, clinical, and replication requirements. It reports no experiment or therapeutic result. Its value is to make future evidence auditable, to prevent surrogate overclaiming, and to define the conditions under which a bounded causal conclusion could be justified.

\appendices
\section{Minimal STABLE-AGE Run Manifest}
\label{app:manifest}

Each future execution records: protocol and consent versions; eligibility and exclusion rules; intervention and comparator definitions; randomization and concealment method; site, visit, and specimen schedule; assay and calibration specifications; endpoint dictionary and direction rules; composite weights; safety monitoring and stop rules; adherence, rescue-care, and discontinuation policy; event definitions; missingness and competing-risk plan; analysis code and environment; data-access conditions; all deviations; all adverse events; all exclusions; null and synthetic-injection tests; replication protocol; and an explicit statement that the run cannot certify permanent cessation of human aging.

\section{Protocol Checklist}
\label{app:checklist}

Before data collection, freeze the construct, population, horizon, intervention boundary, comparator, endpoints, and safety plan. Before unblinding, verify assay identity, replicates, batch controls, duplicate specimens, calibration rules, missingness provenance, and analysis code. During interpretation, report every component domain, adverse event, loss to follow-up, protocol deviation, and exploratory analysis. Before communication, state whether the conclusion is an observation, measurement inference, causal estimate, or long-horizon extrapolation; do not use permanent, universal, or immortalizing language for a finite study.

\begin{thebibliography}{10}

\bibitem{lopez2013}
C. L\'opez-Ot\'in, M. A. Blasco, L. Partridge, M. Serrano, and G. Kroemer, ``The hallmarks of aging,'' \emph{Cell}, vol. 153, no. 6, pp. 1194--1217, 2013, doi: 10.1016/j.cell.2013.05.039.

\bibitem{lopez2023}
C. L\'opez-Ot\'in, M. A. Blasco, L. Partridge, M. Serrano, and G. Kroemer, ``Hallmarks of aging: An expanding universe,'' \emph{Cell}, vol. 186, no. 2, pp. 243--278, 2023, doi: 10.1016/j.cell.2022.11.001.

\bibitem{barbi2018}
E. Barbi \emph{et al.}, ``The plateau of human mortality: Demography of longevity pioneers,'' \emph{Science}, vol. 360, no. 6396, pp. 1459--1461, 2018, doi: 10.1126/science.aat3119.

\bibitem{horvath2013}
S. Horvath, ``DNA methylation age of human tissues and cell types,'' \emph{Genome Biology}, vol. 14, Art. R115, 2013, doi: 10.1186/gb-2013-14-10-r115.

\bibitem{levine2018}
M. E. Levine \emph{et al.}, ``An epigenetic biomarker of aging for lifespan and healthspan,'' \emph{Aging}, vol. 10, no. 4, pp. 573--591, 2018, doi: 10.18632/aging.101414.

\bibitem{kraus2019}
W. E. Kraus \emph{et al.}, ``2 years of calorie restriction and cardiometabolic risk (CALERIE): Exploratory outcomes of a multicentre, phase 2, randomised controlled trial,'' \emph{The Lancet Diabetes \& Endocrinology}, vol. 7, no. 9, pp. 673--683, 2019, doi: 10.1016/S2213-8587(19)30151-2.

\bibitem{belsky2023}
D. W. Belsky \emph{et al.}, ``Caloric restriction slows pace of biological aging in healthy adults,'' \emph{Nature Aging}, vol. 3, pp. 248--257, 2023, doi: 10.1038/s43587-022-00357-y.

\bibitem{mannick2014}
J. B. Mannick \emph{et al.}, ``mTOR inhibition improves immune function in the elderly,'' \emph{Science Translational Medicine}, vol. 6, no. 268, Art. 268ra179, 2014, doi: 10.1126/scitranslmed.3009892.

\bibitem{justice2019}
J. N. Justice \emph{et al.}, ``Senolytics in idiopathic pulmonary fibrosis: Results from a first-in-human, open-label, pilot study,'' \emph{EBioMedicine}, vol. 40, pp. 554--563, 2019, doi: 10.1016/j.ebiom.2018.12.052.

\bibitem{nia2026}
National Institute on Aging, ``What do we know about healthy aging?'' U.S. National Institutes of Health, accessed Sep. 2, 2026. [Online]. Available: \url{https://www.nia.nih.gov/health/healthy-aging/what-do-we-know-about-healthy-aging}

\end{thebibliography}

\end{document}
